% ===========================================================================
%  Bagian 5.5 -- eblup_twofold
% ===========================================================================
\section{Model sub-area dua-lipat: \texttt{eblup\_twofold}}\label{sec:twofold}

\subsection{Setting data dan kapan dipakai}\label{sec:twofold-setting}

Data memiliki tiga tingkat: area (mis.\\ kabupaten), \emph{sub-area} di dalam
area (mis.\\ kecamatan), dan unit sampel di dalam sub-area. Satu baris data =
satu unit sampel, dengan kolom penanda \code{domain} (area) dan
\code{subarea}. Model ini dipakai ketika tersedia kovariat pada tingkat unit
tetapi target estimasi pada dua tingkat sekaligus (area dan sub-area),
termasuk sub-area yang tidak tersampel. Rujukan aslinya
\citet{torabi2014small}; roxygen fungsi ini menyebut model dua-lipat tanpa
mengutip DOI (\code{R/eblup\_twofold.R:7-14}).

\subsection{Persamaan model}\label{sec:twofold-model}

Dua tingkat persamaan (komentar \code{eblup\_twofold.cpp:4-8}):
\begin{align}
  y_{ij} &= \theta_{ij}+e_{ij}, & e_{ij}&\sim N(0,\psi_{ij}),
    \label{eq:tf-lvl1}\\
  \theta_{ij} &= x_{ij}'\beta+v_i+u_{ij}, &
    v_i&\sim N(0,\sigma_v^2),\quad u_{ij}\sim N(0,\sigma_u^2),
    \label{eq:tf-lvl2}
\end{align}
dengan $v_i$ = efek area (notasi kode: \code{s2v}) dan $u_{ij}$ = efek
sub-area (\code{s2u}). Untuk tiap area $d$ kovariasinya adalah
\begin{equation}
  V_d=\sigma_v^2\ones\ones'
  +\mathrm{diag}(\sigma_u^2+\psi_d),
  \qquad
  S_d=\sigma_v^2\ones\ones'+\sigma_u^2\Imat,
  \label{eq:tf-block}
\end{equation}
dan inversinya memakai identitas Sherman--Morrison per blok
\eqref{eq:woodbury-tfh} (biaya $O(n_d)$, bukan $O(n_d^3)$).

\subsection{Prediktor titik}\label{sec:twofold-prediktor}

\begin{equation}
  \hat v_i=\sigma_v^2\,\ones'V^{-1}r,\qquad
  \hat u_{ij}=\sigma_u^2\,(V^{-1}r)_{ij},\qquad
  \hat\theta_{ij}=x_{ij}'\betah+\hat v_i+\hat u_{ij},
  \label{eq:tf-pred}
\end{equation}
dengan $r=y-\Xmat\betah$ dan $V^{-1}$ per blok lewat \eqref{eq:woodbury-tfh}
(baris 351-353). Tiga kasus domain (baris 333-363):
\begin{itemize}
  \item sub-area tersampel pada area tersampel: \eqref{eq:tf-pred};
  \item sub-area \emph{tak}-sampel pada area tersampel:
        $\hat\theta_{ij}=x_{ij}'\betah+\hat v_i$ dan $\hat u=0$ (data hilang
        hanya menyisakan informasi area);
  \item area tanpa sub-area tersampel sama sekali: $\hat\theta=x'\betah$
        (prediktor sintetis).
\end{itemize}
Domain tanpa sampel didefinisikan lewat respons \code{y = NA}
(baris 290-292), sama seperti FH. Validasi mensyaratkan minimal dua area
tersampel dan $N_s>p$ (baris 301, 311).

\subsection{Estimasi komponen varians}\label{sec:twofold-varians}

Fisher scoring dua parameter $(\sigma_v^2,\sigma_u^2)$, nilai awal
$0.5\,\mathrm{median}(\psi)$ keduanya (baris 181-182), toleransi $10^{-4}$,
maksimum 100 iterasi (baris 186). Per iterasi (fungsi \code{twofold\_pass},
baris 80-169) dihitung $Q=(\Xmat'V^{-1}\Xmat)^{-1}$, $r=V^{-1}(y-\Xmat\betah)$,
lalu skor dan informasi:
\begin{align}
  s_v&=-\tfrac12 q_d+\tfrac12(\ones'r)^2,
    & s_u&=-\tfrac12\tr P+\tfrac12 r'r,
    & q_d&=\ones'V^{-1}\ones-(\Xmat'V^{-1}\ones)'Q(\Xmat'V^{-1}\ones),
    \label{eq:tf-score-reml}\\
  \mathcal I_{vv}&=\tfrac12 q_d^2,
    & \mathcal I_{uu}&=\tfrac12\tr(P^2),
    & \mathcal I_{vu}&=\tfrac12\bigl[\ones'V^{-1}\ones
      -2\,\mathbf z'Q\bar x+\bar x'QMQ\bar x\bigr],
  \label{eq:tf-info-reml}
\end{align}
untuk REML (baris 131-151); untuk ML (baris 152-158) suku $\tr P$ dan $q_d$
diganti oleh $\ones'V^{-1}\ones$ dan $\tr V$. Langkah
$\Delta=\mathcal I^{-1}s$ dengan fallback diagonal $s_k/\mathcal I_{kk}$
bila inversi gagal (baris 193-204), nilai negatif dipaksa $0$ (baris 210-211),
dan konvergensi relatif (baris 213-215).

\implnote{Dua hal yang tidak ada pada FH: (i) varian estimasi komponen
diambil dari \emph{invers matriks informasi Fisher} final, bukan bentuk
tertutup $\VarA=2/\sum w^2$ (TFH-3); (ii) loglik REML dihitung benar,
$\ell_R=-\tfrac12[(n-p)\log2\pi+\text{core}+\log|\Xmat'V^{-1}\Xmat|
-\log|\Xmat'\Xmat|]$ (baris 259-265), berbeda dengan FH dan SFH yang selalu
memakai bentuk ML (TFH-9 juga: dasar AIC/BIC memakai jumlah area dengan
sub-area tersampel, bukan jumlah baris).}

\subsection{Estimasi MSE}\label{sec:twofold-mse}

\paragraph{Analitik (default).} Untuk sub-area tersampel (baris 369-411):
\begin{align}
  g_1&=\diag\bigl(S-SV^{-1}S\bigr),
    \label{eq:tf-g1}\\
  g_2&=\diag(DQD'),\qquad D=(\Imat-SV^{-1})\Xmat_d,
    \label{eq:tf-g2}\\
  g_3&=\mathbb E_{vv}\,\mathrm{var}_v+\mathbb E_{uu}\,\mathrm{var}_u
      +2\,\mathbb E_{vu}\,\mathrm{cov}_{vu},
    \label{eq:tf-g3}
\end{align}
di mana $\mathrm{var}_v,\mathrm{var}_u,\mathrm{cov}_{vu}$ adalah elemen
$(\mathcal I^{-1})$ final (baris 324, 407) dan
$\mathbb E_{uu}=\diag(M_uVM_u')$ dengan $M_u=V^{-1}-SV^{-2}$, sedangkan
$M_v=\ones(\ones'V^{-1})'-SV^{-1}\ones(\ones'V^{-1})'$ (baris 395-406).
Totalnya
\begin{equation}
  \widehat{\mathrm{MSE}}=g_1+g_2+g_3,
  \label{eq:tf-mse}
\end{equation}
bukan $g_1+g_2+2g_3$: faktor dua pada suku silang sudah termasuk di dalam
$g_3$ persis seperti tertulis di komentar kode.

\implnote{Di sinilah perbedaan paling menonjol dengan publikasi.
Komentar kode (\code{eblup\_twofold.cpp:19-24}) menyatakan bahwa persamaan
(3.4) \citet{torabi2014small} mengandung suku $[A][B]$ yang semu dan koefisien
varian $u$ yang salah bentuk (plus tanda salah pada lampiran), sehingga
pendekatan matriks $g_3$ dipakai dan divalidasi dengan turunan beda
pemisah (korelasi 0,97/0,96/0,74 terhadap uji numerik). Laporan ini
melaporkannya sebagai TFH-1: apa yang dikodekan adalah pendekatan matriks,
bukan rumus tertutup Torabi--Rao.}

Untuk sub-area tak-sampel (baris 415-468): area tanpa sub-area tersampel
memperoleh $x'Qx+\sigma_v^2+\sigma_u^2$ (baris 421-427); sub-area tak-sampel
pada area tersampel memperoleh $g_1^\ast+g_2^\ast+g_3^\ast$ dengan
$g_1^\ast=\sigma_v^2-\sigma_v^4(\ones'V^{-1}\ones)+\sigma_u^2$ dan
$g_2^\ast=(x-\sigma_v^2\Xmat'V^{-1}\ones)'Q(x-\sigma_v^2\Xmat'V^{-1}\ones)$
(baris 445-464). Nilai negatif dalam rentang $(-10^{-8},0)$ dipangkas ke $0$;
lebih negatif dibiarkan sehingga RSE menjadi \code{NaN} (TFH-8).

\paragraph{Bootstrap parametrik.} Dengan \code{mse = "bootstrap"} (default
$B=200$): semua gambaran dibuat lebih dulu secara sekuensial
($v^\ast\in\mathbb R^{m\times B}$, $u^\ast\in\mathbb R^{N\times B}$,
$e^\ast\in\mathbb R^{N_s\times B}$; baris 475-478), lalu region
\texttt{\#pragma omp parallel for schedule(static)} (baris 481) berjalan per
replikat: $\theta^\ast=\Xmat\betah+v^\ast_i+u^\ast_j$,
$y^\ast=\theta^\ast_{\text{sampel}}+e^\ast$, refit Fisher scoring penuh,
prediktor bootstrap dengan prediktor yang \emph{sama} seperti estimasi titik
(termasuk cabang tak-sampel), dan $\widehat{\mathrm{MSE}}=\mathrm{rata}
\bigl((\hat\theta^\ast-\theta^\ast)^2\bigr)$ (baris 486-517). Karena semua
RNG dipanggil sebelum pragma, hasil tidak bergantung jumlah thread
(TFH-6: pemakaian prediktor yang sama persis juga dinyatakan kode sebagai
perbaikan bug).

\subsection{Pseudo-code}\label{sec:twofold-algo}

\begin{algorithm}[htbp]
\caption{\texttt{eblup\_twofold} --- alur sebenarnya (\texttt{R/eblup\_twofold.R:64-140}, \texttt{src/eblup\_twofold.cpp:173-553})}
\label{alg:tfh}
\KwIn{formula, \code{vardir}, \code{domain} (wajib), \code{subarea},
  \code{method}, \code{mse}, \code{B}, \code{seed}, \code{maxiter}, \code{precision}}
\KwOut{\code{df\_eblup} (area $\times$ sub-area), \code{random\_effect\_var} = $(\sigma_v^2,\sigma_u^2)$, \code{goodness}}
\code{domain} $\leftarrow$ wajib (80-82);\quad
  \code{subarea} $\leftarrow$ \code{seq\_len(nrow(data))} bila \code{NULL} (84-88)\;
\code{area\_idx} $\leftarrow$ \texttt{as.integer(factor(domain)) - 1L} \tcp*{0-based (106)}
\If{\code{seed} tidak \code{NULL}}{\texttt{set.seed(seed)} (108)}\;
$\ldots\leftarrow$ \texttt{.eblup\_twofold\_core($\Xmat,y,\psi$, area, method, mse\_type, B, maxiter, precision)} (110-120)\;
\caption*{}
\KwData{}
\tcp*{\emph{Lapis C++}}
\textbf{gagal} bila \code{method}/\code{mse\_type} tak sah, $N_s\le p$,
  $\psi\le0$, atau area tersampel $<2$ (281-311)\;
$\sigma_v^2\leftarrow\sigma_u^2\leftarrow0.5\,\mathrm{median}(\psi)$ (181-182)\;
\While{\texttt{diff} $>$ \code{precision} \textbf{dan} $k<$ \code{maxiter}}{
  per blok $d$: $V_d^{-1}a\leftarrow$ \eqref{eq:woodbury-tfh} (98-108, 345-350)\;
  $\betah,r,q_d,\tr P,\tr V,\tr P^2\leftarrow$ akumulasi (110-145)\;
  \uIf{ML}{$s,\mathcal I\leftarrow$ bentuk ML (152-158)}
  \uElse{$s,\mathcal I\leftarrow$ \eqref{eq:tf-score-reml}--\eqref{eq:tf-info-reml} (131-151)}\;
  $\Delta\leftarrow\mathcal I^{-1}s$; bila inversi gagal $\to\Delta_k=s_k/\mathcal I_{kk}$ (193-204)\;
  $\sigma_k^2\leftarrow\max(\sigma_k^2+\Delta_k,0)$; non-finite $\to$ nilai lama (206-211)\;
  $\texttt{diff}\leftarrow\max_k\abs{(\sigma_k^{2(new)}-\sigma_k^2)/\max(\sigma_k^2,10^{-12})}$ (213-215)\;
}
pass final pada varian final: $\betah,Q,\mathcal I^{-1}$ (223-237)\;
$\hat\theta\leftarrow$ \eqref{eq:tf-pred}, termasuk cabang tak-sampel (333-363)\;
\uIf{\code{mse} $=$ analytical}{$g_1,g_2,g_3\leftarrow$ \eqref{eq:tf-g1}--\eqref{eq:tf-g3};
  $\mathrm{MSE}\leftarrow$ \eqref{eq:tf-mse} per baris tersampel (369-411);
  cabang tak-sampel (415-468)}
\uElse{$v^\ast,u^\ast,e^\ast\leftarrow$ bangkit sekuensial (475-478)\;
  \texttt{\#pragma omp parallel for schedule(static)} (481)\;
  \For{$b=0,\dots,B-1$ paralel}{$\theta^\ast,y^\ast\leftarrow$ (486-488);
    refit \code{twofold\_fit} (490-491); $\mathrm{SqDiff}[:,b]\leftarrow
    (\hat\theta^\ast-\theta^\ast)^2$ (495-517)}}\;
$\ell,\mathrm{AIC},\mathrm{BIC}\leftarrow$ loglik ML/REML, dasar $m_{\text{fit}}$ (239-267, 545-548)\;
susun \code{df\_eblup} (\code{domain, subarea, y, eblup, vardir, random\_effect\_area,
random\_effect\_subarea, mse, rse}) (550-553)\;
\end{algorithm}

\subsection{Signature, argumen, objek keluaran, dan contoh}\label{sec:twofold-api}

\begin{lstlisting}[language=R,caption={Signature \texttt{eblup\_twofold} (\texttt{R/eblup\_twofold.R:64-77}).}]
eblup_twofold(formula, vardir, domain = NULL, subarea = NULL, data,
              method = c("REML", "ML"),
              mse = c("analytical", "bootstrap"),
              B = 200, seed = NULL, maxiter = 100,
              precision = 1e-4, print_result = TRUE)
\end{lstlisting}

\begin{itemize}
  \item \code{domain} --- wajib (bila \code{NULL} fungsi berhenti);
        \code{subarea} --- opsional, default penomoran unit.
  \item \code{mse} --- \code{"analytical"} (default) atau \code{"bootstrap"}
        ($B$ default 200).
  \item \code{seed} --- dipanggil pada semua jalur, bukan hanya bootstrap
        (TFH-7).
\end{itemize}

Objek keluaran: \code{"fastsae"} berisi \code{random\_effect\_var} bernama
\code{sigma2\_v} (area) dan \code{sigma2\_u} (sub-area), \code{estcoef},
\code{df\_eblup}, \code{goodness} (\code{loglikelihood, AIC, BIC} dengan
$\mathrm{AIC}=-2\ell+2(p+2)$ dan $\mathrm{BIC}=-2\ell+(p+2)\log m_{\text{fit}}$),
\code{n\_iter}, \code{convergence}, \code{method = "eblup"},
\code{level = "subarea"}, \code{model = "TWOFOLD"}.

\begin{lstlisting}[language=R,caption={Contoh yang dapat dijalankan (dari \texttt{man/eblup\_twofold.Rd}); membutuhkan kolom penanda area dan sub-area.}]
library(fastsae)
m1 <- eblup_twofold(y ~ x1 + x2, data = your_data, vardir = "vardir",
                    domain = "area", subarea = "kec")
m1$random_effect_var        # sigma2_v (area), sigma2_u (subarea)

m2 <- eblup_twofold(y ~ x1 + x2, data = your_data, vardir = "vardir",
                    domain = "area", subarea = "kec",
                    mse = "bootstrap", B = 100, seed = 42)
\end{lstlisting}

\subsection{Referensi kunci}\label{sec:twofold-ref}

Rujukan model dua-lipat: \citet{torabi2014small} (disebut
\code{DESCRIPTION} beserta DOI-nya dan menjadi dasar komentar kode
\code{eblup\_twofold.cpp:19-24}); kerangka EBLUP: \citet{rao2015small};
pembanding: \pkg{sae} \citep{molina2015sae}.
